\documentclass[10pt,a4paper]{article}

% --- Margins (Maximizes printable area safely) ---
\usepackage[top=0.5cm, bottom=0.7cm, left=0.5cm, right=0.7cm]{geometry}

% --- Multi-column Layout ---
\usepackage{multicol}
\setlength{\columnsep}{0.4cm} 
\setlength{\columnseprule}{0.2pt}

% --- Math & Fonts ---
\usepackage{amsmath, amssymb, amsfonts}
\DeclareMathSizes{10}{10}{7}{6} % {fontsize}{textstyle}{scriptstyle}{scriptscriptstyle}
\usepackage{bm}
\usepackage{array}
\usepackage{booktabs}
\usepackage{xcolor}

% --- Hebrew & English Support ---
\usepackage{polyglossia}
\setmainlanguage{hebrew}
\setotherlanguage{english}
\newfontfamily\hebrewfont{David CLM}[Script=Hebrew, Scale=1.1]

\usepackage{bidi}

\usepackage[inline]{enumitem}
% --- Tight Lists ---
\usepackage{enumitem}
\setlist{
    noitemsep,                   % מרווח אפסי בין איטמים
    topsep=1pt,                  % מרווח אפסי מעל/מתחת לרשימה
    partopsep=0pt,
    parsep=0pt,
    leftmargin=6pt              % צמצום ההזחה הצידית של הרשימה
}
% \setlist{nolistsep, leftmargin=*} 


% --- Compact Section Headings ---
\usepackage{titlesec}
% \titleformat{\section}{\large\bfseries}{\thesection}{0.4em}{}
% \titlespacing*{\section}{0pt}{3pt}{1pt}

% פורמט כותרת \section קומפקטית
\titleformat{\section}
  {\normalfont\bfseries\small\color{blue!80!black}}
  {\thesection}{1em}{}
\titlespacing*{\section}
  {0pt}{1pt}{1pt} % {left}{before-space}{after-space}



% --- Compact Math Spacing ---
\AtBeginDocument{
  \setlength{\abovedisplayskip}{1pt}
  \setlength{\belowdisplayskip}{1pt}
  \setlength{\abovedisplayshortskip}{1pt}
  \setlength{\belowdisplayshortskip}{1pt}
}

\pagestyle{empty} % Removes page numbers to save vertical space
\raggedbottom

\usepackage{tabularx}
\usepackage{array}


\begin{document}
\fontsize{10pt}{12pt}\selectfont

\begin{multicols*}{2}

\section{פורמליזם יחסותי, וקטורים, טנזורים ואינדקסים}
המטריקה ב \textenglish{Mostly Minus} היא 
$
\eta_{\mu\nu} = \eta^{\mu\nu} = \text{diag}(1, -1, -1, -1)
$
\begin{itemize}
    \item \textbf{4-וקטור קונטרה (אינדקס עליון):} $A^\mu = (A^0, \vec{A})$. משתנה תחת $A'^\mu = \Lambda^\mu{}_\nu A^\nu = \dfrac{\partial x'^\mu}{\partial x^\nu} A^\nu \implies \mathbf{A}' = \mathbf{J} \mathbf{A}$.
    \item \textbf{4-וקטור קו (אינדקס תחתון):} $A_\mu=(A_0, -\vec{A})$. משתנה תחת $A'_\mu =\Lambda_\mu{}^\nu A_\nu = (\Lambda^{-1})^\nu{}_\mu A_\nu = \dfrac{\partial x^\nu}{\partial x'^\mu} A_\nu \Rightarrow \mathbf{A}' = (\mathbf{J}^{-1})^T \mathbf{A}$.
    \item \textbf{העלאת/הורדת אינדקסים:} $A_\mu = \eta_{\mu\nu} A^\nu$, $A^\mu = \eta^{\mu\nu} A_\nu$.
    \item \textbf{טרנספורמציה של טנזורים (כתיב אינדקסים):} 
    \[
    T'^{\mu\nu} = \Lambda^\mu{}_\alpha \Lambda^\nu{}_\beta T^{\alpha\beta}, \
    T'^\mu{}_\nu = \Lambda^\mu{}_\alpha \Lambda_\nu{}^\beta T^\alpha{}_\beta, 
    T'_{\mu\nu} = \Lambda_\mu{}^\alpha \Lambda_\nu{}^\beta T_{\alpha\beta}
    \]

    \item \textbf{טרנספורמציה של טנזורים (כתיב מטריצי):}
    \begin{itemize}
        \item \textbf{עליונים ($T^{\mu\nu}$):} $\mathbf{T}' = \Lambda \mathbf{T} \Lambda^T \quad | \quad \mathbf{T} = \Lambda^{-1} \mathbf{T}' (\Lambda^{-1})^T$
        \item \textbf{תחתונים ($T_{\mu\nu}$):} $\mathbf{T}' = (\Lambda^{-1})^T \mathbf{T} \Lambda^{-1} \quad | \quad \mathbf{T} = \Lambda^T \mathbf{T}' \Lambda$
        \item \textbf{מעורבים ($T^\mu{}_\nu$):} $\mathbf{T}' = \Lambda \mathbf{T} \Lambda^{-1} \quad | \quad \mathbf{T} = \Lambda^{-1} \mathbf{T}' \Lambda$
    \end{itemize}

\item \textbf{זמן עצמי ($\tau$) והאינטרוואל השמור ($ds$):}
\begin{itemize}
    \item \textbf{קשר לאורך השמור:} $ds^2 = c^2 d\tau^2 = \eta_{\mu\nu} dx^\mu dx^\nu$
    \item \textbf{קשר לזמן במערכת כללית ($t$):} $d\tau = \frac{ds}{c} = \frac{dt}{\gamma}$
    \item \textbf{אינטגרל במסלול:} $\tau = \int \frac{ds}{c} = \int_{t_1}^{t_2} \sqrt{1 - \frac{v^2(t)}{c^2}} \, dt$
\end{itemize}

\item \textbf{ייצוג טנזורים כמטריצות :}
$
T^\mu{}_\nu \quad \leftrightarrow \quad T_{\text{שורה}, \, \text{עמודה}} = T_{\mu, \nu}
$

\item \textbf{טרנספוז:}
$
    (F^T)^{\mu\nu} = F^{\nu\mu}, \ (g^T)_{\mu\nu} = g_{\nu\mu}, \ (T^T)^\mu{}_\nu = T_\nu{}^\mu
$

\end{itemize}

\section{איזומטריה (טרנספורמציית סימטריה)}
טרנספ' $x'^\mu = \Lambda^\mu{}_\nu x^\nu$ היא איזומטריה של מ' מינקובסקי אם היא שומרת על המטריקה:
$
\eta_{\alpha\beta} \Lambda^\alpha{}_\mu \Lambda^\beta{}_\nu = \eta_{\mu\nu} \Leftrightarrow \Lambda^T \eta \Lambda = \eta
$

\section{4-וקטורים חשובים}
\begin{itemize}
    \item \textbf{מיקום:} $x^\mu = (ct, \vec{r})$, \quad $x_\mu = (ct, -\vec{r})$
    \item \textbf{נגזרת:} $\partial_\mu \equiv \frac{\partial}{\partial x^\mu} = \left(\frac{1}{c}\frac{\partial}{\partial t}, \vec{\nabla}\right)$, \ $\partial^\mu \equiv \frac{\partial}{\partial x_\mu} = \left(\frac{1}{c}\frac{\partial}{\partial t}, -\vec{\nabla}\right)$
    \item \textbf{4-מהירות:} $U^\mu = \frac{dx^\mu}{d\tau} = \gamma(c, \vec{v})$, כאשר $U^\mu U_\mu = c^2$
    \item \textbf{4-תאוצה:} $A^\mu = \frac{dU^\mu}{d\tau} = \gamma \frac{d}{dt}\left[\gamma(c, \vec{v})\right] = \gamma\left(c \dot{\gamma}, \, \dot{\gamma}\vec{v} + \gamma\vec{a}\right)$
    \\
\begin{itemize*}[label=-, itemjoin={,\quad}]
    \item נגזרת $\gamma$: $\dot{\gamma} = \frac{d\gamma}{dt} = \frac{\gamma^3}{c^2}(\vec{v} \cdot \vec{a})$
    \item א"ג ל-4-מהירות: $U_\mu A^\mu = 0$
\end{itemize*}
    \item \textbf{4 תנע-אנ':} $P^\mu = m U^\mu = \left(\frac{E}{c}, \vec{p}\right) = \gamma m(c, \vec{v})$, עם $P^\mu P_\mu = m^2 c^2$
    \item \textbf{4-כח (לורנץ):} $K^\mu = \frac{dP^\mu}{d\tau} = m A^\mu = \frac{q}{c} F^{\mu\nu} U_\nu = \gamma \left( \frac{\vec{F} \cdot \vec{v}}{c}, \vec{F} \right)$. \\ {\small$\vec{F} = q (\vec{E} + \frac{\vec{v}}{c} \times \vec{B}) = \frac{d\vec{p}}{dt} = \frac{d(\gamma m \vec{v})}{dt} | \frac{\vec{p}}{E} = \frac{\gamma m\vec{v}}{\gamma m c^2} \Rightarrow \vec{v} = \frac{\vec{p}c^2}{E} | \frac{\vec{F} \cdot \vec{v}}{c}=\frac{dE}{d(ct)}$}
    \item \textbf{4-פוטנציאל:} $A^\mu = (\Phi, \vec{A})$, \quad $A_\mu = (\Phi, -\vec{A})$
    \item \textbf{4 צפיפות זרם:} $J_\mu = (\rho c, -\vec{J}), \ J^\mu = \rho_0 u^\mu = (\rho c, \vec{J})$ \\ ($\text{ צפיפות במע' מנוחה} \rho_0, \; \rho = \gamma \rho_0$)
\end{itemize}

\section{משוואות מקסוול וטנזור השדה $F_{\mu\nu}$}
\begin{itemize}
    \item 
    \begin{itemize*}[label=, itemjoin={\quad | \quad}]
        \item \textbf{הגדרה:} $F_{\mu\nu} = \partial_\mu A_\nu - \partial_\nu A_\mu$
        \item $F^{ij} = \epsilon^{ijk} B^k$
    \end{itemize*}
    \item 
    {\scriptsize
    \setlength{\arraycolsep}{1.5pt}
    $
    F_{\mu\nu} = \begin{pmatrix}
    0 & E_x & E_y & E_z \\
    -E_x & 0 & -B_z & B_y \\
    -E_y & B_z & 0 & -B_x \\
    -E_z & -B_y & B_x & 0
    \end{pmatrix}, \;
    F^{\mu\nu} = \begin{pmatrix}
    0 & -E_x & -E_y & -E_z \\
    E_x & 0 & -B_z & B_y \\
    E_y & B_z & 0 & -B_x \\
    E_z & -B_y & B_x & 0
    \end{pmatrix}
    $
    }
    \item 
    $
    E^i = F^{i0} = -F^{0i} = F_{0i} = -F_{i0}, \:
    B^i = \frac{1}{2}\epsilon^{ijk} F^{jk} = \frac{1}{2}\epsilon^{ijk} F_{jk}
    $
    \item \textbf{משוואות מקסוול (אי-הומוגניות):}
    \[
    \partial_\mu F^{\mu\nu} = \frac{4\pi}{c} J^\nu \iff \vec{\nabla}\cdot\vec{E} = 4\pi\rho, \quad \vec{\nabla}\times\vec{B} - \frac{1}{c}\frac{\partial\vec{E}}{\partial t} = \frac{4\pi}{c}\vec{J}
    \]
    \item \textbf{זהות ביאנקי (משוואות הומוגניות):}
    \[
    \partial_\alpha F_{\beta\gamma} + \partial_\beta F_{\gamma\alpha} + \partial_\gamma F_{\alpha\beta} = 0 \iff \partial_\mu \tilde{F}^{\mu\nu} = 0
    \]
    \[
    \vec{\nabla}\cdot\vec{B} = 0, \quad \vec{\nabla}\times\vec{E} + \frac{1}{c}\frac{\partial\vec{B}}{\partial t} = 0
    \]
\end{itemize}

\textbf{B,E באמצעות הפוט' $(\Phi, \vec{A})$:}
$
\vec{E} = -\vec{\nabla}\Phi - \frac{1}{c}\frac{\partial \vec{A}}{\partial t}, \ \vec{B} = \vec{\nabla} \times \vec{A}
$

\textbf{תנאי שפה וקפיצה על משטח ($\sigma$):} 
\\
$\Delta E_\perp = 4\pi\sigma, \; \Delta \vec{E}_\parallel = 0, \; \Delta\Phi = 0, \; \left.\dfrac{\partial\Phi_2}{\partial n} - \dfrac{\partial\Phi_1}{\partial n}\right|_{\text{surface}} = -4\pi\sigma$

\section{הטנזור הדואלי $\tilde{F}^{\mu\nu}$}
מוגדר ע"י $\tilde{F}^{\mu\nu} = \frac{1}{2}\epsilon^{\mu\nu\alpha\beta}F_{\alpha\beta}$. מתקבל מהחלפת השדות: $\vec{E} \to \vec{B}$ ו-$\vec{B} \to -\vec{E}$.
משמש לכתיבת המשוואות ההומ' בקיצור: $\partial_\mu \tilde{F}^{\mu\nu} = 0$.

\section{גדלים אינווריאנטיים של השדה}
\begin{itemize*}[label=\textbullet, itemjoin={\quad}]
    \item $F_{\mu\nu}F^{\mu\nu} = 2(B^2 - E^2)$
    \item $F_{\mu\nu}\tilde{F}^{\mu\nu} = -4(\vec{E}\cdot\vec{B})$
\end{itemize*}

\section{טנזור מאמץ-אנרגיה $T_{\mu\nu}^{A}$}
\begin{itemize}
    \item \textbf{גדלים אלקטרומגנטיים \textenglish{(CGS)}:}
    \begin{itemize}
        \item צפיפות אנרגיה: $u = \frac{1}{8\pi}(E^2+B^2)$
        \item שטף אנרגיה (וקטור פוינטינג): $\vec{S} = \frac{c}{4\pi}(\vec{E}\times\vec{B})$
        \item צפיפות תנע: $\vec{p}_{\text{EM}} = \dfrac{1}{c^2}\vec{S} = \dfrac{1}{4\pi c}(\vec{E}\times\vec{B})$
        \item צפיפות תנע זוויתי: $\vec{\mathcal{L}}_{\text{EM}} = \vec{r} \times \vec{p}_{\text{EM}} = \frac{1}{c^2}(\vec{r} \times \vec{S})$
        \item טנזור המאמצים: $\sigma_{ij} = \frac{1}{4\pi}\left(E_i E_j + B_i B_j - \frac{1}{2}\delta_{ij}(E^2+B^2)\right)$
    \end{itemize}

    \item 
    $T^{\mu\nu} = -\frac{1}{4\pi c} \left( F^\mu{}_\alpha F^{\nu\alpha} - \frac{1}{4}\eta^{\mu\nu}F_{\alpha\beta}F^{\alpha\beta} \right) 
    = \begin{pmatrix} 
    \frac{u}{c} & \vec{p}_{\text{EM}}^T \\
    \frac{1}{c^2}\vec{S} & -\frac{1}{c}\overleftrightarrow{\sigma}
    \end{pmatrix}
    $

    \item $
    T_{\mu\nu} = -\frac{1}{4\pi c} \left( F_{\mu\alpha}F_\nu{}^\alpha - \frac{1}{4}\eta_{\mu\nu}F_{\alpha\beta}F^{\alpha\beta} \right)
    = \begin{pmatrix} 
    \frac{u}{c} & -\vec{p}_{\text{EM}}^T \\
    -\frac{1}{c^2}\vec{S} & -\frac{1}{c}\overleftrightarrow{\sigma}
    \end{pmatrix}
    $

    \item $T^\mu{}_\nu = -\frac{1}{4\pi c} \left( F^{\mu\alpha} F_{\nu\alpha} - \frac{1}{4}\delta^\mu_\nu F_{\alpha\beta}F^{\alpha\beta} \right) 
    = \begin{pmatrix} 
    \frac{u}{c} & -\vec{p}_{\text{EM}}^T \\
    \frac{1}{c^2}\vec{S} & \frac{1}{c}\overleftrightarrow{\sigma}
    \end{pmatrix}
    $

    \item \textbf{דיברגנץ טנזור מאמץ-אנ' (שימור תנע-אנ' וצפיפות 4-כוח לורנץ):}
    \[
    \partial_\mu T^\mu{}_\nu = -\frac{1}{c^2} F_{\nu\rho} J^\rho, \quad \partial_\mu T^{\mu\nu} = -\frac{1}{c^2} F^{\nu\rho} J_\rho
    \]

    \begin{itemize}
        \item \textbf{רכיב הזמן ($\nu=0$ - שימור אנרגיה):} $\frac{\partial u}{\partial t} + \vec{\nabla} \cdot \vec{S} = -\vec{J} \cdot \vec{E}$
        \item \textbf{' מרחביים ($\nu=i$ - '' תנע):} $\frac{\partial \vec{p}_{\text{EM}}}{\partial t} - \vec{\nabla} \cdot \overleftrightarrow{\sigma} = -\left( \rho \vec{E} + \frac{1}{c} \vec{J} \times \vec{B} \right)$
        \item \textbf{מע' סגורה (ללא מטענים/זרמים, $J_\mu=0$):} $\partial_\mu T^\mu{}_\nu = 0$
        \item 
        $
        (\nabla \cdot \mathbf{\sigma})_i = \partial_j \sigma_{ji} = \frac{\partial \sigma_{ji}}{\partial x_j}
        $
    \end{itemize}
    \item \textbf{עקבת הטנזור \textenglish{(Trace)}:} $\text{Tr}(T) = T^\mu{}_\mu = 0$.
\end{itemize}

\section*{סוגי כיול (\textenglish{Gauge Transformations}) - טרנספ' כיול: $A_\mu \to A_\mu + \partial_\mu \Lambda$}
\begin{itemize}
    \item \textbf{כיול לורנץ:} תנאי: 
    $
    \partial_\mu A^\mu = 0 \iff \frac{1}{c}\frac{\partial \Phi}{\partial t} + \vec{\nabla}\cdot\vec{A} = 0
    $
    מתאים לבעיות יחסותיות וגלים א"מ. מעביר את משוואות מקסוול לצורה:
    \[
    \square A^\mu = \frac{4\pi}{c}J^\mu \quad \left(\square = \partial_\alpha \partial^\alpha = \frac{1}{c^2}\frac{\partial^2}{\partial t^2} - \nabla^2\right)
    \]
    \item \textbf{כיול קולון (\textenglish{Coulomb Gauge}):} תנאי: $\vec{\nabla}\cdot\vec{A} = 0$.\\
    מתאים לא"ס, מ"ס ולקוונטיזציה של השדה. נותן:
    \[
    \nabla^2 \Phi = -4\pi\rho \implies \Phi(\vec{r}, t) = \int \frac{\rho(\vec{r}', t)}{|\vec{r}-\vec{r}'|} d^3 r'
    \]
    \item \textbf{כיול טמפורלי:} ($ \phi =A^0 = 0$). השדה החשמלי - $\vec{E} = -\frac{1}{c}\frac{\partial \vec{A}}{\partial t}$.

\item \textbf{כיול אקסיאלי:} קובעים רכיב מרחבי: $A_z = 0$ (בכללי $\vec{n} \cdot \vec{A} = 0$).
\end{itemize}

\section{הערך האינטגרלי (הפעולה) ומשוואות שדה}

\begin{itemize}
    \item \textbf{הפעולה הא"מ ($S$):}
    $
    S = \int \mathcal{L}_{\text{EM}} \, d^4x = \int \left( \mathcal{L}_{\text{field}} + \mathcal{L}_{\text{int}} \right) d^4x
    $
    כאשר $d^4x = c \, dt \, d^3x$.

    \item \textbf{צפיפות הלגראנז'יאן:}
    $
    \mathcal{L}_{\text{EM}} = -\frac{1}{16\pi c} F_{\mu\nu}F^{\mu\nu} - \frac{1}{c^2} J_\mu A^\mu
    $
    \begin{itemize}
    \item \textbf{איבר השדה החופשי:} $\mathcal{L}_{\text{field}} = -\frac{1}{16\pi c} F_{\mu\nu}F^{\mu\nu} = \frac{1}{8\pi c} \left( E^2 - B^2 \right)$
        \item \textbf{איבר האינטראקציה:} $\mathcal{L}_{\text{int}} = -\frac{1}{c^2} J_\mu A^\mu = -\frac{1}{c} \rho \Phi + \frac{1}{c^2} \vec{J} \cdot \vec{A}$
    \end{itemize}

    \item \textbf{משוואות \textenglish{EL} לשדה ($A_\nu$):}
    $
    \partial_\mu \left( \dfrac{\partial \mathcal{L}_{\text{EM}}}{\partial (\partial_\mu A_\nu)} \right) - \dfrac{\partial \mathcal{L}_{\text{EM}}}{\partial A_\nu} = 0
    $

    \item \textbf{משוואת \textenglish{EL} לשדה סקלרי ($\phi$):}
$
\partial_\mu \left( \dfrac{\partial \mathcal{L}}{\partial (\partial_\mu \phi)} \right) - \dfrac{\partial \mathcal{L}}{\partial \phi} = 0
$

    \item \textbf{נגזרות:}
    $
    \dfrac{\partial \mathcal{L}_{\text{EM}}}{\partial (\partial_\mu A_\nu)} = -\dfrac{1}{4\pi c} F^{\mu\nu}, \; \dfrac{\partial \mathcal{L}_{\text{EM}}}{\partial A_\nu} = -\dfrac{1}{c^2} J^\nu
    $

    \item \textbf{$\delta S = 0$ $\Leftarrow$ מ' מקסוול לא-הומוגניות:}
    $
    \partial_\mu F^{\mu\nu} = \frac{4\pi}{c} J^\nu
    $
\end{itemize}

\section{טרנספורמציות לורנץ וסיבובים}
\begin{itemize}
    \item \textbf{בוסט ב$x$:} \textenglish{Rapidity} ($\theta=\text{tanh}^{-1}(\beta)$), $\cosh\theta=\gamma, \sinh\theta=\gamma\beta$
    {\scriptsize
    \setlength{\arraycolsep}{2.5pt}
    \[
    \Lambda^\mu{}_\nu = \begin{pmatrix}
    \gamma & -\gamma\beta & 0 & 0 \\
    -\gamma\beta & \gamma & 0 & 0 \\
    0 & 0 & 1 & 0 \\
    0 & 0 & 0 & 1
    \end{pmatrix} = \begin{pmatrix}
    \cosh\theta & -\sinh\theta & 0 & 0 \\
    -\sinh\theta & \cosh\theta & 0 & 0 \\
    0 & 0 & 1 & 0 \\
    0 & 0 & 0 & 1
    \end{pmatrix}, \ \gamma=\frac{1}{\sqrt{1-\beta^2}}
    , \beta=\frac{v}{c}\]
    }
    \item \textbf{סיבוב סביב $z$ ($x'^\mu = \Lambda_z(\theta)^\mu{}_\nu x^\nu$):}
    {\scriptsize
    \setlength{\arraycolsep}{2.5pt}
    $
    \Lambda_z(\theta) = \begin{pmatrix}
    1 & 0 & 0 & 0 \\
    0 & \cos\theta & \sin\theta & 0 \\
    0 & -\sin\theta & \cos\theta & 0 \\
    0 & 0 & 0 & 1
    \end{pmatrix}
    $
    }
\end{itemize}

\section*{טנזור לוי-צ'יוויטה ( $E_{\alpha\beta\gamma\delta} = \text{sgn}(g)\sqrt{|g|}\,\epsilon_{\alpha\beta\gamma\delta}, \ E^{\alpha\beta\gamma\delta} = \frac{1}{\sqrt{|g|}}\,\epsilon^{\alpha\beta\gamma\delta}$)}
\begin{itemize}
     \item \textbf{3 אינ' (מרחביים):} $\epsilon_{123} = -\epsilon^{123} = +1$. אנטי-סימ' להחלפת אינ'.
    \item \textbf{4 ':} במינקובסקי (\textenglish{Mostly Minus}): $\epsilon^{0123} = +1$, לכן $\epsilon_{0123} = -1$.
\end{itemize}

\section{אופרטורים דיפרנציאליים ($g \equiv \det(g_{\mu\nu})$)}

\begin{itemize}
    \item \textbf{גרדיאנט:} 
    $ \vec{\nabla}\phi = \frac{\partial \phi}{\partial x^i} e^i = \partial_i \phi \, e^i
    $
    ($\partial_i \phi$ קו-וואריאנטי.)

    \item \textbf{דיברגנס:} עבור שדה וקטורי
    $
    \vec{\nabla} \cdot \vec{E} = \frac{1}{\sqrt{g}} \partial_j \left( \sqrt{g} \, E^j \right)
    $

    \item \textbf{רוטור:} הרכיב הקונטרה 
    $ (\vec{\nabla} \times \vec{E})^i = \frac{\varepsilon^{ijk}}{\sqrt{g}} \frac{\partial E_k}{\partial x^j} = \frac{\varepsilon^{ijk}}{\sqrt{g}} \partial_j E_k
    $ \\
    \textbf{מעבר לרכיב פיזיקלי (בכיוון $\hat{e}_i$):} $A^{(\hat{i})} = \sqrt{g_{ii}} \, A^i = \frac{1}{\sqrt{g_{ii}}} \, A_i$

    \item \textbf{לפלסיאן:}
    \[
    \nabla^2 \phi = \vec{\nabla} \cdot \vec{\nabla}\phi = \frac{1}{\sqrt{g}} \frac{\partial}{\partial x^j} \left( \sqrt{g} \, g^{jk} \frac{\partial \phi}{\partial x^k} \right) = \frac{1}{\sqrt{g}} \partial_j \left( \sqrt{g} \, g^{jk} \partial_k \phi \right)
    \]
    שימו לב כי $\partial^i(\cdot) \neq \partial_i(g^{ji}\cdot)$.
    הנגזרת עם אינדקס למעלה מוגדרת $\partial^i(\cdot) = g^{ji} \partial_j(\cdot)$. כי הקוא' עקומות

    \item \textbf{קשר בין וקטורי יחידה $\hat{e}_\mu$ לווקטורי בסיס מטבעיים $e_\mu = \frac{\partial}{\partial x^\mu}$:}
\begin{itemize}
    \item \textbf{גליליות:} $\hat{\rho} = e_\rho = e^\rho, \hat{\phi} = \frac{1}{\rho} e_\phi = \rho e^\phi, \hat{z} = e_z = e^z, \left(e_\phi = \rho \hat{\phi}\right)$
    \item \textbf{ספ':} $\hat{r} = e_r, \ \hat{\theta} = \dfrac{1}{r} e_\theta, \ \hat{\phi} = \dfrac{1}{r \sin\theta} e_\phi, \ \left(e_\theta = r\hat{\theta}, \; e_\phi = r\sin\theta \hat{\phi}\right)$
\end{itemize}
\end{itemize}


\section{פונקציות גרין ופתרון משוואת פואסון}
\begin{itemize}
    \item \textbf{חישוב הפוטנציאל הפואסוני:} $\nabla^2 \Phi = -4\pi \rho$
    \item \textbf{הגדרת פונקציית גרין:} $\nabla'^2 G(\vec{r}, \vec{r}') = -4\pi \delta^{(3)}(\vec{r} - \vec{r}')$
    \item \textbf{כללי-סינגולרי+הרמוני:}
    {\small$
    G(\vec{r}, \vec{r}') = \frac{1}{|\vec{r} - \vec{r}'|} + F(\vec{r}, \vec{r}'), \nabla'^2 F(\vec{r}, \vec{r}') = 0
    $}
    \item \textbf{סימטריה:} $G(\vec{r}, \vec{r}') = G(\vec{r}', \vec{r})$
    \item \textbf{נוסחת הפתרון הכללית (משפט גרין):}
    \[
    \Phi(\vec{r}) = \int_V \rho(\vec{r}') G(\vec{r}, \vec{r}') \, d^3r' - \frac{1}{4\pi} \oint_S \left[ \Phi(\vec{r}') \frac{\partial G}{\partial n'} - G \frac{\partial \Phi}{\partial n'} \right] dA'
    \]
    \item \textbf{תנאי דיריכלה ($G_D = 0$ על $S$):}
    \[
    \Phi(\vec{r}) = \int_V \rho(\vec{r}') G_D(\vec{r}, \vec{r}') \, d^3r' - \frac{1}{4\pi} \oint_S \Phi(\vec{r}') \frac{\partial G_D}{\partial n'} \, dA'
    \]
    \item \textbf{תנאי נוימן ($\frac{\partial G_N}{\partial n'} = -\frac{4\pi}{S}$ על $S$):}
    \[
    \Phi(\vec{r}) = \langle \Phi \rangle_S + \int_V \rho(\vec{r}') G_N(\vec{r}, \vec{r}') \, d^3r' + \frac{1}{4\pi} \oint_S G_N(\vec{r}, \vec{r}') \frac{\partial \Phi}{\partial n'} \, dA'
    \]
\end{itemize}

\section{פונקציות גרין ומטעני דמות}
\begin{itemize}
    \item \textbf{הגדרת $r_<$ ו-$r_>$:} $r_< = \min(r, r')$, \quad $r_> = \max(r, r')$.

    \item \textbf{מישור אינסופי מוארק ($z=0$):} 
    מטען מקורי $q$ ממוקם ב-$\vec{r}' = (x', y', z')$ (כאשר $z' > 0$). מטען דמות $q' = -q$ ב-$\vec{r}_{\text{img}} = (x', y', -z')$ ($z' < 0$):
    $
    G(\vec{r},\vec{r}') = \dfrac{1}{|\vec{r}-\vec{r}'|} - \dfrac{1}{|\vec{r}-\vec{r}_{\text{img}}|}
    $
    

    \item \textbf{ספירה מוארקת ברדיוס $R$ (תנאי דיריכלה):}
    מטען $q$ ממוקם ב-$\vec{r}'$. מטען דמות $q' = -q \frac{R}{r'}$ ממוקם ב-$\vec{r}_{\text{img}} = \frac{R^2}{r'^2}\vec{r}'$.
    \[
    G(\vec{r},\vec{r}') = \frac{1}{|\vec{r}-\vec{r}'|} - \frac{R/r'}{\left|\vec{r} - [(R^2)/(r'^2)]\vec{r}'\right|}
    \]
    כאשר $\cos\gamma = \hat{r}\cdot\hat{r}' = \cos\theta\cos\theta' + \sin\theta\sin\theta'\cos(\phi-\phi')$:
    \[
    G(\vec{r},\vec{r}') = \frac{1}{\sqrt{r^2 + r'^2 - 2rr'\cos\gamma}} - \frac{1}{\sqrt{(rr'/R)^2 + R^2 - 2rr'\cos\gamma}}
    \]

    \item \textbf{פונקציית גרין דיריכלה לספירה ברדיוס $R$ (פיתוח ספרי):}
$$ G_D(\vec{r},\vec{r}') = \sum_{\ell=0}^\infty \sum_{m=-\ell}^\ell \frac{4\pi}{2\ell+1} \left[ \frac{r_<^\ell}{r_>^{\ell+1}} - \mathcal{R}_\ell \right] Y_{\ell m}^*(\theta',\phi') Y_{\ell m}(\theta,\phi) $$
כאשר $\mathcal{R}_\ell = (rr')^\ell / R^{2\ell+1}$ עבור $r,r' < R$ (פנים), \\ ו-$\mathcal{R}_\ell = R^{2\ell+1} / (rr')^{\ell+1}$ עבור $r,r' > R$ (חוץ).
\end{itemize}

\section{תכונות זוגיות (Parity) של פולינומי לז'נדר והרמוניות ספריות:}
\begin{itemize}
    \item $P_\ell(-x) = (-1)^\ell P_\ell(x) \Rightarrow P_\ell(\cos(\pi - \theta)) = (-1)^\ell P_\ell(\cos\theta)$
    \item \textbf{פולינומים נלווים:} $P_\ell^m(-x) = (-1)^{\ell-m} P_\ell^m(x)$
    \item \textbf{שיקוף מרחבי ($\vec{r} \to -\vec{r}$):} $\theta \to \pi - \theta, \; \phi \to \phi + \pi \Rightarrow Y_{\ell,m}(\pi - \theta, \, \phi + \pi) = (-1)^\ell Y_{\ell,m}(\theta, \phi)$
    \item \textbf{צמוד וזווית $\phi$:} $Y_{\ell, -m}(\theta, \phi) = (-1)^m Y_{\ell, m}^*(\theta, \phi), \ Y_{\ell,m}(\theta, -\phi) = (-1)^m Y_{\ell,-m}(\theta, \phi)$
\end{itemize}

\section{פוטנציאל בסימטריה כדורית / פיתוח מולטיפולי}
פוטנציאל מפיזור מטען $\rho(\vec{r}')$:
\[
\Phi(\vec{r}) = 4\pi \sum_{\ell=0}^\infty \sum_{m=-\ell}^\ell \frac{q_{\ell,m}}{2\ell+1} \frac{Y_{\ell,m}(\theta,\phi)}{r^{\ell+1}} \quad (r > r_{\max} \text{עבור })
\]
כאשר המומנטים המולטיפוליים הם:
\[
q_{\ell,m} = \int \rho(\vec{r}') r'^\ell Y_{\ell,m}^*(\theta',\phi') d^3r', \quad q_{\ell,-m} = (-1)^m q_{\ell,m}^*
\]

\textbf{פוטנציאל של קליפה כדורית ברדיוס $R$ - בעיית לפלס:}
\begin{itemize}
    \item \textbf{כללי (תלוי ב-$r,\theta,\phi$, במונחי הרמוניות כדוריות $Y_{\ell,m}$):}
    \[
\phi(r, \theta, \phi) = \sum_{\ell=0}^{\infty} \sum_{m=-\ell}^{\ell} \left( A_{\ell,m} r^\ell + B_{\ell,m} r^{-(\ell+1)} \right) Y_{\ell,m}(\theta, \phi)
\]
    
    \item \textbf{סימ' אזימ' ($m=0$):} מצטמצם לפ' לז'נדר $P_\ell(\cos\theta)$ עם $A_\ell, B_\ell$.

    \item \textbf{תכונות אורתוגונליות ונירמול:}
    \begin{itemize}
        \item \textbf{פולינומי לז'נדר ($P_\ell(x)$):} 
         \[
        \int_{-1}^{1} P_\ell(x) P_{\ell'}(x) \, dx = \int_{0}^{\pi} P_\ell(\cos\theta) P_{\ell'}(\cos\theta) \sin\theta \, d\theta = \frac{2}{2\ell+1} \delta_{\ell\ell'}
        \]
        \item
        $
        \int_{-1}^{1} P_\ell^m(x) P_{\ell'}^m(x) \, dx = \frac{2}{2\ell+1} \frac{(\ell+m)!}{(\ell-m)!} \delta_{\ell\ell'}
        $
        \item $
        \int_{0}^{2\pi} d\phi \int_{0}^{\pi} \sin\theta \, d\theta \, Y_{\ell,m}^*(\theta,\phi) Y_{\ell',m'}(\theta,\phi) = \delta_{\ell\ell'} \delta_{mm'}
        $

        \item \textbf{יחס שלמות:} ($\frac{\delta(\theta - \theta') \delta(\phi - \phi')}{\sin\theta} = \delta(\cos\theta - \cos\theta') \delta(\phi - \phi')$) \\
        $
        \sum\limits_{\ell=0}^{\infty} \sum\limits_{m=-\ell}^{+\ell} Y_{\ell m}^*(\theta, \phi) Y_{\ell m}(\theta', \phi') = \frac{\delta(\theta - \theta') \delta(\phi - \phi')}{\sin\theta}
        $
        
        \item $
    \dfrac{1}{|\vec{r} - \vec{r}'|} = \sum_{l=0}^{\infty} \dfrac{r_<^l}{r_>^{l+1}} P_l(\cos\gamma)$
      \item $
      P_l(\cos\gamma) = \dfrac{4\pi}{2l+1} \sum_{m=-l}^{l} Y_{lm}^*(\theta', \phi') Y_{lm}(\theta, \phi)
      $
    \end{itemize}
\end{itemize}

\section*{פיתוח מולטיפולי קרטזי ($r \gg r'$)}
\noindent
פיתוח טיילור של הפוטנציאל האלקטרוסטטי במרחקים גדולים:
\[
\phi(\vec{r}) = \frac{q}{r} + \frac{\vec{d} \cdot \vec{r}}{r^3} + \sum_{i,j} Q_{ij} \frac{x_i x_j}{r^5} = \dfrac{q}{r} + \dfrac{d_i x_i}{r^3} + Q_{ij} \dfrac{x_i x_j}{r^5}
\]
\noindent
\textbf{הגדרת איברי המולטיפול:}
\begin{itemize}
    \item \textbf{מונופול (מטען כולל):} $q = \int \rho(\vec{r}') \, d^3r'$
    \item \textbf{מומנט דיפול חשמלי:} $\vec{d} = \int \vec{r}' \rho(\vec{r}') \, d^3r'$
    \item \textbf{קוואדרופול ($\text{Tr}(Q)=0$):} 
    $
    Q_{ij} = \frac{1}{2} \int \rho(\vec{r}') \left( 3x'_i x'_j - r'^2 \delta_{ij} \right) d^3r'
    $
\end{itemize}

\section*{פוטנציאל וקטורי מגנטי $\vec{A}$}
\noindent
\textbf{חישוב מגנטוסטטי מאינטגרל צפיפות זרם (כיול קולון $\nabla \cdot \vec{A} = 0$):}
\noindent
{\small\[\vec{A}(\vec{r}) = \frac{1}{c} \int \frac{\vec{J}(\vec{r}')}{|\vec{r} - \vec{r}'|} \, d^3r', \ \vec{B}(\vec{r}) = \nabla \times \vec{A} = \frac{1}{c} \int \frac{\vec{J}(\vec{r}') \times (\vec{r} - \vec{r}')}{|\vec{r} - \vec{r}'|^3} \, d^3r'\]}
\begin{itemize}
    \item \textbf{איבר הדיפול המגנטי ($n=1$):} האיבר המוביל במרחקים גדולים:
    \[
    \vec{A}_{\text{dip}}(\vec{r}) = \frac{\vec{m} \times \vec{r}}{r^3} = \frac{\vec{m} \times \hat{r}}{r^2}, \ \vec{B}(\vec{r}) = \nabla \times \vec{A}_{\text{dip}}(\vec{r}) = \frac{3(\vec{m} \cdot \hat{r})\hat{r} - \vec{m}}{r^3}
    \]
    \item \textbf{מומנט דיפול מגנטי $\vec{m}$:}
$
\vec{m} = \frac{1}{2c} \int (\vec{r}' \times \vec{J}(\vec{r}')) \, d^3r'
$
\end{itemize}

\section*{קרינה מולטיפולית והספק (תלוי בזמן, כיול לורנץ $\nabla \cdot \vec{A} + \frac{1}{c}\frac{\partial \phi}{\partial t} = 0$)}
\noindent
{\small
$
\vec{A}_{\text{rad}}(\vec{r},t) = \dfrac{1}{cr} \iiint d^3r' \, \vec{J}\left(\vec{r}', t - \dfrac{r}{c} + \dfrac{\hat{r} \cdot \vec{r}'}{c}\right)
$
}\\
{\small$
\vec{B}_{\text{rad}}(\vec{r},t) = -\dfrac{\hat{r}}{c} \times \dfrac{\partial \vec{A}_{\text{rad}}}{\partial t},
\vec{E}_{\text{rad}}(\vec{r},t) = \hat{r} \times \left[ \dfrac{\hat{r}}{c} \times \dfrac{\partial \vec{A}_{\text{rad}}}{\partial t} \right] = -\hat{r} \times \vec{B}_{\text{rad}}(\vec{r},t)
$}
\noindent
\textbf{נוסחת לֶרְמוֹר (לא-יחסותי $v \ll c$):}
$P = \dfrac{2}{3} \dfrac{q^2 a^2}{c^3} \ \left(a = |\ddot{\vec{r}}_0|\right)$

\noindent
\textbf{חישוב התאוצה $a$ ללרמור:} $\vec{a} = \frac{\vec{F}}{m} = \frac{q}{m}\left(\vec{E} + \frac{\vec{v}}{c}\times\vec{B}\right)$ \quad $|$ \quad \textbf{תנועה מעגלית:} $a = \frac{v^2}{R} = \omega^2 R$ \quad $|$ \quad \textbf{מתנד:} $\langle a^2 \rangle = \frac{1}{2}\omega^4 x_0^2$

\section{וקטורי יחידה במערכות קואורדינטות}
\begin{itemize}

    \item \textbf{כדוריות לקרטזיות:}
    $\hat{r} = \sin\theta\cos\phi\hat{x} + \sin\theta\sin\phi\hat{y} + \cos\theta\hat{z}$, \quad
    $\hat{\theta} = \cos\theta\cos\phi\hat{x} + \cos\theta\sin\phi\hat{y} - \sin\theta\hat{z}$, \quad
    $\hat{\phi} = -\sin\phi\hat{x} + \cos\phi\hat{y}$
    \item \textbf{קרטזיות לכדוריות:}
    $\hat{x} = \sin\theta\cos\phi\hat{r} + \cos\theta\cos\phi\hat{\theta} - \sin\phi\hat{\phi}$, \quad
    $\hat{y} = \sin\theta\sin\phi\hat{r} + \cos\theta\sin\phi\hat{\theta} + \cos\phi\hat{\phi}$, \quad
    $\hat{z} = \cos\theta\hat{r} - \sin\theta\hat{\theta}$
    \item \textbf{גליליות לק':}
    $\hat{r} = \cos\phi\hat{x} + \sin\phi\hat{y}$, \
    $\hat{\phi} = -\sin\phi\hat{x} + \cos\phi\hat{y}$, \
    $\hat{z} = \hat{z}$
    \item \textbf{קרט' לגליליות:}
    $\hat{x} = \cos\phi\hat{r} - \sin\phi\hat{\phi}$, \
    $\hat{y} = \sin\phi\hat{r} + \cos\phi\hat{\phi}$, \
    $\hat{z} = \hat{z}$

    \item \textbf{כדוריות ($\hat{r}, \hat{\theta}, \hat{\phi}$):}
    $\hat{r} \times \hat{\theta} = \hat{\phi}, \quad \hat{\theta} \times \hat{\phi} = \hat{r}, \quad \hat{\phi} \times \hat{r} = \hat{\theta}$
    \item \textbf{גליליות ($\hat{r}, \hat{\phi}, \hat{z}$):}
    $\hat{r} \times \hat{\phi} = \hat{z}, \quad \hat{\phi} \times \hat{z} = \hat{r}, \quad \hat{z} \times \hat{r} = \hat{\phi}$

\end{itemize}

\section*{קרינת מולטיפולים (שדה רחוק $r \gg \lambda \gg d$, כאשר $\lambda = \frac{2\pi c}{\omega}$, יחידות \textenglish{CGS})}

\noindent
\textbf{קירובים:} 
{\small $r \gg d$: \textbf{מקור נקודתי} ($\dfrac{1}{|\vec{r}-\vec{r}'|} \approx \frac{1}{r}$) \ $|$ \ 
$r \gg \lambda$: \textbf{אזור קרינה} (איברי $r^{-1}$, $\vec{E} \perp \vec{B} \perp \hat{r}$) \ $|$ \ 
$\lambda \gg d \Leftrightarrow \dfrac{v}{c} \ll 1, \ (v \approx d \cdot \omega)$: \textbf{לא-יחסותי} (הפרדה למולטיפולים).
\textbf{קיצורים וסימונים:} $t' = t - r/c$, \quad $\vec{E} = \vec{B} \times \hat{r}$, \quad $\vec{U}_{\mathbf{Q}} \equiv \dddot{\mathbf{Q}} \cdot \hat{r}$, $\quad U_r \equiv \hat{r} \cdot \dddot{\mathbf{Q}} \cdot \hat{r}$. כאשר $(\mathbf{Q} \cdot \hat{r})_i = \sum_j Q_{ij} \hat{r}_j$}. \textbf{חתך פעולה דיפ':} $\frac{d\sigma}{d\Omega} = \frac{1}{|\vec{S}_{in}|} \left\langle \frac{dP_{out}}{d\Omega} \right\rangle$.
\vspace{3pt}
\begingroup
% 1. הגדלת מקדם גובה השורות מ-1.9 ל-2.3 (מרווח נדיב שמונע חיתוך נוסחאות)
\renewcommand{\arraystretch}{2.2}
\footnotesize
\setlength{\tabcolsep}{2pt}

\begin{tabularx}{\linewidth}{|c|c|*{3}{>{\centering\arraybackslash}X|}}
\hline
\textbf{גודל} & \textbf{מונופול} ($q$) & \textbf{ED} ($\vec{p}$) & \textbf{MD} ($\vec{m}$) & \textbf{EQ} ($Q_{ij}$) \\ \hline

$\Phi(\vec{r},t)$ 
& $\dfrac{q}{r}$ 
& $\dfrac{\dot{\vec{p}} \cdot \hat{r}}{cr}$ 
& $0$ 
& $\dfrac{\hat{r} \cdot \ddot{\mathbf{Q}} \cdot \hat{r}}{3c^2 r}$ \\ \hline

$\vec{A}(\vec{r},t)$ 
& $0$ 
& $\dfrac{\dot{\vec{p}}}{cr}$ 
& $\dfrac{\dot{\vec{m}} \times \hat{r}}{cr}$ 
& $\dfrac{\ddot{\mathbf{Q}} \cdot \hat{r}}{3c^2 r}$ \\ \hline

$\vec{B}(\vec{r},t)$ 
& $0$ 
& $\dfrac{\ddot{\vec{p}} \times \hat{r}}{c^2 r}$ 
& $\dfrac{(\ddot{\vec{m}} \cdot \hat{r})\hat{r} - \ddot{\vec{m}}}{c^2 r}$ 
& $-\dfrac{\hat{r} \times \vec{U}_{\mathbf{Q}}}{3c^3 r}$ \\ \hline

$\vec{E}(\vec{r},t)$ 
& $\dfrac{q\hat{r}}{r^2}$ 
& $\dfrac{(\ddot{\vec{p}} \times \hat{r}) \times \hat{r}}{c^2 r}$ 
& $\dfrac{\hat{r} \times \ddot{\vec{m}}}{c^2 r}$ 
& $-\dfrac{\vec{U}_{\mathbf{Q}} - U_r\hat{r}}{3c^3 r}$ \\ \hline

$\dfrac{dP}{d\Omega} = r^2 \hat{r} \cdot \vec{S}$ 
& $0$ 
& $\dfrac{|\ddot{\vec{p}} \times \hat{r}|^2}{4\pi c^3}$ 
& $\dfrac{|\hat{r} \times \ddot{\vec{m}}|^2}{4\pi c^3}$ 
& $\dfrac{|\vec{U}_{\mathbf{Q}}|^2 - |U_r|^2}{36\pi c^5}$ \\ \hline

$P_{\text{total}}$ 
& $0$ 
& $\dfrac{2}{3c^3}|\ddot{\vec{p}}|^2$ 
& $\dfrac{2}{3c^3}|\ddot{\vec{m}}|^2$ 
& $\dfrac{1}{45 c^5} \sum\limits_{i,j} |\dddot{Q}_{ij}|^2$ \\ \hline
\end{tabularx}
\endgroup

\section{אופרטורים וקטוריים בקואורדינטות עקומות אורתוגונליות}
\begin{itemize}
\item \textbf{גליליות ($\rho, \phi, z$):}
\begin{itemize}[leftmargin=*]
    \item \textbf{גרדיאנט:} $\nabla f = \frac{\partial f}{\partial \rho}\hat{\rho} + \frac{1}{\rho}\frac{\partial f}{\partial \phi}\hat{\phi} + \frac{\partial f}{\partial z}\hat{z}$
    \item \textbf{דיברגנץ:} $\nabla \cdot \vec{A} = \frac{1}{\rho}\frac{\partial(\rho A_\rho)}{\partial \rho} + \frac{1}{\rho}\frac{\partial A_\phi}{\partial \phi} + \frac{\partial A_z}{\partial z}$
    \item \textbf{רוטור:} $\nabla \times \vec{A} = \left(\frac{1}{\rho}\frac{\partial A_z}{\partial \phi} - \frac{\partial A_\phi}{\partial z}\right)\hat{\rho} + \left(\frac{\partial A_\rho}{\partial z} - \frac{\partial A_z}{\partial \rho}\right)\hat{\phi} + \frac{1}{\rho}\left(\frac{\partial(\rho A_\phi)}{\partial \rho} - \frac{\partial A_\rho}{\partial \phi}\right)\hat{z}$
    \item \textbf{לפלסיאן:} $\nabla^2 f = \frac{1}{\rho}\frac{\partial}{\partial \rho}\left(\rho \frac{\partial f}{\partial \rho}\right) + \frac{1}{\rho^2}\frac{\partial^2 f}{\partial \phi^2} + \frac{\partial^2 f}{\partial z^2}$
\end{itemize}

\item \textbf{כדוריות ($r, \theta, \phi$):}
\begin{itemize}[leftmargin=*]
    \item \textbf{גרדיאנט:} $\nabla f = \frac{\partial f}{\partial r}\hat{r} + \frac{1}{r}\frac{\partial f}{\partial \theta}\hat{\theta} + \frac{1}{r\sin\theta}\frac{\partial f}{\partial \phi}\hat{\phi}$
    \item \textbf{דיברגנץ:} $\nabla \cdot \vec{A} = \frac{1}{r^2}\frac{\partial(r^2 A_r)}{\partial r} + \frac{1}{r\sin\theta}\frac{\partial(\sin\theta A_\theta)}{\partial \theta} + \frac{1}{r\sin\theta}\frac{\partial A_\phi}{\partial \phi}$
    \item \textbf{רוטור:} $\nabla \times \vec{A} = \frac{1}{r\sin\theta}\left(\frac{\partial(\sin\theta A_\phi)}{\partial \theta} - \frac{\partial A_\theta}{\partial \phi}\right)\hat{r} + \frac{1}{r}\left(\frac{1}{\sin\theta}\frac{\partial A_r}{\partial \phi} - \frac{\partial(r A_\phi)}{\partial r}\right)\hat{\theta} + \frac{1}{r}\left(\frac{\partial(r A_\theta)}{\partial r} - \frac{\partial A_r}{\partial \theta}\right)\hat{\phi}$
    \item \textbf{לפלסיאן:} $\nabla^2 f = \frac{1}{r^2}\frac{\partial}{\partial r}\left(r^2 \frac{\partial f}{\partial r}\right) + \frac{1}{r^2\sin\theta}\frac{\partial}{\partial \theta}\left(\sin\theta \frac{\partial f}{\partial \theta}\right) + \frac{1}{r^2\sin^2\theta}\frac{\partial^2 f}{\partial \phi^2}$
\end{itemize}
\end{itemize}

\section{מעבר מטריצי בין בסיסים אורתונורמליים ($R^{-1} = R^T$)}
\begin{itemize}

\item \textbf{גליליות $\leftrightarrow$ קרטזיות:}
{\small$$\begin{pmatrix} \hat{\rho} \\ \hat{\phi} \\ \hat{z} \end{pmatrix} = 
\begin{pmatrix} 
\cos\phi & \sin\phi & 0 \\ 
-\sin\phi & \cos\phi & 0 \\ 
0 & 0 & 1 
\end{pmatrix} 
\begin{pmatrix} \hat{x} \\ \hat{y} \\ \hat{z} \end{pmatrix}
 \mid \begin{pmatrix} \hat{x} \\ \hat{y} \\ \hat{z} \end{pmatrix} = 
\begin{pmatrix} 
\cos\phi & -\sin\phi & 0 \\ 
\sin\phi & \cos\phi & 0 \\ 
0 & 0 & 1 
\end{pmatrix} 
\begin{pmatrix} \hat{\rho} \\ \hat{\phi} \\ \hat{z} \end{pmatrix}$$}

\item \textbf{כדוריות $\leftrightarrow$ קרטזיות:}
$$\begin{pmatrix} \hat{r} \\ \hat{\theta} \\ \hat{\phi} \end{pmatrix} = 
\begin{pmatrix} 
\sin\theta\cos\phi & \sin\theta\sin\phi & \cos\theta \\ 
\cos\theta\cos\phi & \cos\theta\sin\phi & -\sin\theta \\ 
-\sin\phi & \cos\phi & 0 
\end{pmatrix} 
\begin{pmatrix} \hat{x} \\ \hat{y} \\ \hat{z} \end{pmatrix}$$

$$\begin{pmatrix} \hat{x} \\ \hat{y} \\ \hat{z} \end{pmatrix} = 
\begin{pmatrix} 
\sin\theta\cos\phi & \cos\theta\cos\phi & -\sin\phi \\ 
\sin\theta\sin\phi & \cos\theta\sin\phi & \cos\phi \\ 
\cos\theta & -\sin\theta & 0 
\end{pmatrix} 
\begin{pmatrix} \hat{r} \\ \hat{\theta} \\ \hat{\phi} \end{pmatrix}$$
\end{itemize}

\section{פונקציות גרין, ייצוגים בקואורדינטות ונגזרות שפה}

\begin{itemize}[leftmargin=*]

    \item \textbf{1. מישור אינסופי מוארק ($z'=0$, נפח הבעיה: $z > 0$):}
    \begin{itemize}
        \item \textbf{ייצוג גלילי:}
        $G_{\text{cyl}} = \frac{1}{\sqrt{\rho^2 + \rho'^2 - 2\rho\rho'\cos(\Delta\phi) + (z-z')^2}} - \frac{1}{\sqrt{\rho^2 + \rho'^2 - 2\rho\rho'\cos(\Delta\phi) + (z+z')^2}}$
        \item \textbf{נגזרת מכוונת על השפה ($\hat{n}' = -\hat{z} \implies \frac{\partial G}{\partial n'} = -\frac{\partial G}{\partial z'}$):}
        \begin{itemize}
            \item \textbf{קרטזית:} $\left.\frac{\partial G}{\partial n'}\right|_{z'=0} = \frac{2z}{\left[(x-x')^2 + (y-y')^2 + z^2\right]^{3/2}}$
            \item \textbf{גלילית:} $\left.\frac{\partial G}{\partial n'}\right|_{z'=0} = \frac{2z}{\left[\rho^2 + \rho'^2 - 2\rho\rho'\cos(\phi-\phi') + z^2\right]^{3/2}}$
            \item \textbf{ספרית ($\theta'=\frac{\pi}{2}$):} $\left.\frac{\partial G}{\partial n'}\right|_{\theta'=\pi/2} = \frac{2r\cos\theta}{\left[r^2 + r'^2 - 2rr'\sin\theta\cos(\phi-\phi')\right]^{3/2}}$
        \end{itemize}
    \end{itemize}

    \item \textbf{2. ספירה מוארקת ברדיוס $R$ (תנאי דיריכלה, שפה: $r'=R$):}
    \begin{itemize}
        \item \textbf{פונקציית גרין (ספרית):} $\left(q' = -q\frac{R}{r'}, \; \vec{r}_{\text{img}} = \frac{R^2}{r'^2}\vec{r}'\right)$
        \item \textbf{ייצוג גלילי:}
        $G_{\text{cyl}} = \frac{1}{\sqrt{\rho^2+\rho'^2-2\rho\rho'\cos(\Delta\phi)+(z-z')^2}} - \frac{1}{\sqrt{\frac{(\rho^2+z^2)(\rho'^2+z'^2)}{R^2} + R^2 - 2(\rho\rho'\cos(\Delta\phi)+zz')}}$
        \item \textbf{נגזרת מכוונת על השפה ($\frac{\partial G}{\partial n'} = \mp \frac{\partial G}{\partial r'}$ עבור בעיה חיצונית/פנימית):}
        \begin{itemize}
            \item \textbf{ספרית:} $\left.\frac{\partial G}{\partial n'}\right|_{r'=R} = \pm \frac{r^2 - R^2}{R\left(r^2 + R^2 - 2rR\cos\gamma\right)^{3/2}}$
            \item \textbf{קרטזית:} $\left.\frac{\partial G}{\partial n'}\right|_{r'=R} = \pm \frac{x^2 + y^2 + z^2 - R^2}{R\left[x^2 + y^2 + z^2 + R^2 - 2(xx' + yy' + zz')\right]^{3/2}}$
            \item \textbf{גלילית:} $\left.\frac{\partial G}{\partial n'}\right|_{r'=R} = \pm \frac{\rho^2 + z^2 - R^2}{R\left[\rho^2 + z^2 + R^2 - 2R(\rho\sin\theta'\cos\Delta\phi + z\cos\theta')\right]^{3/2}}$
        \end{itemize}
        *(כאשר $+$ מיועד לבעיה חיצונית $r \ge R$, ו-$-$ לבעיה פנימית $r \le R$)*
    \end{itemize}

    \item \textbf{קשר בין נגזרת מכוונת ($\frac{\partial G}{\partial n'}$) לנגזרות קרטזיות ($\nabla' G \cdot \hat{n}'$):}
\begin{itemize}
    \item \textbf{שפה גלילית ($\hat{n}' = \hat{\rho}'$):} 
    $\frac{\partial G}{\partial n'} = \frac{\partial G}{\partial \rho'} = \cos\phi' \frac{\partial G}{\partial x'} + \sin\phi' \frac{\partial G}{\partial y'}$
    
    \item \textbf{שפה כדורית ($\hat{n}' = \hat{r}'$):} 
    $\frac{\partial G}{\partial n'} = \frac{\partial G}{\partial r'} = \sin\theta'\cos\phi' \frac{\partial G}{\partial x'} + \sin\theta'\sin\phi' \frac{\partial G}{\partial y'} + \cos\theta' \frac{\partial G}{\partial z'}$
    
    \item \textbf{שפה מישורית ($\hat{n}' = \hat{z}'$):} 
    $\frac{\partial G}{\partial n'} = \frac{\partial G}{\partial z'}$

\item \textbf{קשר במישור אינסופי ($z'=0 \iff \theta'=\frac{\pi}{2}$) לפי כיוון הנורמל:}
\begin{itemize}
    \item {\small\textbf{$z'>0$ ($\hat{n}' = -\hat{z}'$):}} 
    $\left.\frac{\partial G}{\partial n'}\right|_{z'=0} = -\left.\frac{\partial G}{\partial z'}\right|_{z'=0} = +\frac{1}{r'} \left.\frac{\partial G}{\partial \theta'}\right|_{\theta'=\pi/2}$
    \item \textbf{נורמל בכיוון $+\hat{z}'$:} 
    $\left.\frac{\partial G}{\partial z'}\right|_{z'=0} = -\frac{1}{r'} \left.\frac{\partial G}{\partial \theta'}\right|_{\theta'=\pi/2}$
\end{itemize}
\end{itemize}
\end{itemize}

\section{פונקציית הדלתא של דיראק וזהויות טנזוריות}

\begin{itemize}[leftmargin=*]
    \item \textbf{דלתא של דיראק בקואורדינטות עקומות אורתוגונליות במרחב $D$-ממדי:}
    $$\delta^{(D)}(\vec{r} - \vec{r}') = \frac{1}{\prod_{i=1}^{D} h_i} \prod_{i=1}^{D} \delta(q_i - q_i') = \frac{1}{\sqrt{\det(g)}} \prod_{i=1}^{D} \delta(q_i - q_i')$$
    כאשר $h_i$ הם גורמי קנה-המידה \textenglish{(scale factors)}, $g_{ij}$ הוא טנזור המטריקה, ו-$\det(g) = \prod_{i=1}^{D} h_i^2$ הוא הדטרמיננטה שלו.
    \begin{itemize}
        \item \textbf{גליליות ($\rho, \phi, z$):} $\delta^{(3)}(\vec{r} - \vec{r}') = \frac{1}{\rho} \delta(\rho - \rho') \delta(\phi - \phi') \delta(z - z')$
        \item \textbf{כדוריות ($r, \theta, \phi$):} $\delta^{(3)}(\vec{r} - \vec{r}') = \frac{1}{r^2 \sin\theta} \delta(r - r') \delta(\theta - \theta') \delta(\phi - \phi') = \frac{1}{r^2} \delta(r - r') \delta(\cos\theta - \cos\theta') \delta(\phi - \phi')$
    \end{itemize}

    \item \textbf{זהויות בסיסיות של פונקציית הדלתא:}
    \begin{itemize}
        \item \textbf{שינוי משתנה פונקציונלי ($f(x_i)=0$ שורשים פשוטים):} $\delta(f(x)) = \sum_{i} \frac{\delta(x - x_i)}{|f'(x_i)|}$
        \item \textbf{זהות המרחק הרדיאלי (לפלסיאן של $1/r$):} $\nabla^2 \left( \frac{1}{|\vec{r} - \vec{r}'|} \right) = -4\pi \delta^{(3)}(\vec{r} - \vec{r}')$
        \item \textbf{זהות הנגזרת והסינון:} $\int f(x) \delta'(x - a) \, dx = -f'(a), \quad x \delta'(x) = -\delta(x)$
        \item \textbf{נגזרת פונקציית המדרגה (Heaviside):} $\frac{d}{dx}\Theta(x - a) = \delta(x - a)$
    \end{itemize}

    \item \textbf{הקשר בין לוי-צ'יוויטה ($\epsilon_{ijk}$), קרונקר ($\delta_{ij}$) ופונקציית הדלתא:}
    \begin{itemize}
        \item \textbf{זהות הכיווץ של לוי-צ'יוויטה (מכפלת סביבונים):}
        $$\epsilon_{ijk} \epsilon_{imn} = \delta_{jm}\delta_{kn} - \delta_{jn}\delta_{km} \mid \epsilon_{ijk} \epsilon_{ijn} = 2\delta_{kn}, \quad \epsilon_{ijk} \epsilon_{ijk} = 6$$
        \item \textbf{זהות הרוטור כפול (תחליף לאופרטורים וקטוריים):}\\
        $(\nabla \times (\nabla \times \vec{A}))_i = \epsilon_{ijk} \partial_j (\epsilon_{klm} \partial_l A_m) = \nabla_i (\nabla \cdot \vec{A}) - \nabla^2 A_i$
        \item \textbf{ייצוג נפח רציף מרוכז ע"י דלתא ולוי-צ'יוויטה (צפיפות זרם של מטען נקודתי $q$):}
        $$\rho(\vec{r}) = q \delta^{(3)}(\vec{r} - \vec{r}_0(t)), \quad \vec{J}(\vec{r}) = q \vec{v} \delta^{(3)}(\vec{r} - \vec{r}_0(t))$$
        \text{רכיב הזרם המגנטי/צפיפות: } $$J_i(\vec{r}) = \epsilon_{ijk} m_j \partial_k \delta^{(3)}(\vec{r} - \vec{r}_0) \quad (\vec{m} \text{עבור דיפול מגנטי })$$
    \end{itemize}
\end{itemize}

\section{סימטריות מרחביות של מומנטים מולטיפוליים (חשמליים ומגנטיים)}

\begin{itemize}[leftmargin=*]
    \item \textbf{הגדרת המומנטים והתנהגותם תחת שיקוף מרחבי \textenglish{(Parity $\mathcal{P}$)}:}
    \begin{itemize}
        \item \textbf{מונופול חשמלי ($q$):} סקלר זוגי, $\mathcal{P} q = +q$.
        \item \textbf{דיפול חשמלי ($\vec{p} = \int \rho(\vec{r}') \vec{r}' d^3r'$):} וקטור פולרי (אמיתי), $\mathcal{P} \vec{p} = -\vec{p}$.
        \item \textbf{קוודרופול חשמלי ($Q_{ij} = \int \rho(\vec{r}') [3x'_i x'_j - r'^2 \delta_{ij}] d^3r'$):} טנזור דרוג 2 סימטרי וחסר עקבה, זוגי לשיקוף ($\mathcal{P} Q_{ij} = +Q_{ij}$).
        \item \textbf{דיפול מגנטי ($\vec{m} = \frac{1}{2c} \int \vec{r}' \times \vec{J}(\vec{r}') d^3r'$):} אקסיאל-וקטור (וקטור מדומה), זוגי לשיקוף מרחבי: $\mathcal{P} \vec{m} = +\vec{m}$.
    \end{itemize}

    \item \textbf{תנאי סימטריה המאפסים מומנטים מולטיפוליים:}

    \begin{itemize}
        \item \textbf{סימטריה לשיקוף במרכז ($r' \to -r'$ או \textenglish{Parity}):}
        \begin{itemize}
            \item אם צפיפות המטען אי-זוגית $\rho(-\vec{r}') = -\rho(\vec{r}')$: \textbf{$q = 0$ וגם $Q_{ij} = 0$} (כל המומנטים הזוגיים מתאפסים).
            \item אם צפיפות המטען זוגית $\rho(-\vec{r}') = +\rho(\vec{r}')$: \textbf{$\vec{p} = 0$} (כל המומנטים האי-זוגיים מתאפסים).
            \item אם צפיפות הזרם אי-זוגית $\vec{J}(-\vec{r}') = -\vec{J}(\vec{r}')$: \textbf{$\vec{m} = 0$} (דיפול מגנטי מתאפס).
        \end{itemize}

        \item \textbf{סימטריה לשיקוף במישור (למשל $z \to -z$):}
        \begin{itemize}
            \item אם $\rho(x,y,-z) = -\rho(x,y,z)$: $\Rightarrow$ הרכיב $p_z = 0$, וכן $Q_{xx}, Q_{yy}, Q_{zz}, Q_{xy} = 0$ (רק רכיבי הצימוד $Q_{xz}, Q_{yz}$ יכולים לשרוד).
            \item אם לזרם יש סימטריית שיקוף במישור $xy$: הרכיבים המקבילים למישור של הדיפול המגנטי מתאפסים ($m_x = m_y = 0$).
        \end{itemize}

        \item \textbf{סימטריה סיבובית גלילית \textenglish{(Axial Symmetry)} סביב ציר $z$:}
        \begin{itemize}
            \item תלות בזווית הפלנרית: $\rho(\vec{r}') = \rho(r', \theta')$.
            \item \textbf{דיפול חשמלי:} מתווה אך ורק לאורך ציר הסימטריה: $\vec{p} = p_z \hat{z}$ (כלומר $p_x = p_y = 0$).
            \item \textbf{קוודרופול חשמלי:} הטנזור באלכסון בלבד עם רכיב יחיד בלתי-תלוי:
            $Q_{xx} = Q_{yy} = -\frac{1}{2} Q_{zz}, \quad Q_{ij} = 0 \;\; (i \neq j)$
            \item \textbf{דיפול מגנטי:} לזרם שזורם במעגלים סביב ציר $z$ ($\vec{J} = J_\phi \hat{\phi}$): $\vec{m} = m_z \hat{z}$ ($m_x = m_y = 0$).
        \end{itemize}

        \item \textbf{סימטריה כדורית מלאה \textenglish{(Spherical Symmetry)}:}
        \begin{itemize}
            \item תלות רדיאלית בלבד: $\rho(\vec{r}') = \rho(r')$.
            \item $\Leftarrow$ \textbf{$\vec{p} = 0$ וגם $Q_{ij} = 0$} (כל המומנטים החשמליים מסדר $l \ge 1$ מתאפסים; נשאר רק המונופול $q$).
            \item $\Leftarrow$ \textbf{$\vec{m} = 0$} (כל המומנטים המגנטיים מתאפסים).
        \end{itemize}
    \end{itemize}
\end{itemize}

\section*{בסיס טבעי ($e_\mu = \partial_\mu$), דואלי ($e^\mu = dx^\mu$) ומטריקה מרחבית ($g_{ij}, g^{ij}$):}
\begin{itemize}
    \item \textbf{גליליות:} 
    $\hat{\rho} = e_\rho = e^\rho, \quad \hat{\phi} = \frac{1}{\rho} e_\phi = \rho e^\phi, \quad \hat{z} = e_z = e^z$ \\
    $g_{ij} = \text{diag}(1, \rho^2, 1), \quad g^{ij} = \text{diag}(1, \rho^{-2}, 1)$
    
    \item \textbf{ספ':} 
    $\hat{r} = e_r = e^r, \quad \hat{\theta} = \frac{1}{r} e_\theta = r e^\theta, \quad \hat{\phi} = \frac{1}{r \sin\theta} e_\phi = r\sin\theta e^\phi$ \\
    $g_{ij} = \text{diag}(1, r^2, r^2\sin^2\theta), \quad g^{ij} = \text{diag}(1, r^{-2}, r^{-2}\sin^{-2}\theta)$
\end{itemize}

\section{חילוץ מקדמים בטורים של פונקציות אורתוגונליות}

\begin{itemize}[leftmargin=*]
    \item \textbf{הגדרה כללית של המערכת האורתוגונלית:}
    נתונה פונקציה $f(x)$ המוצגת כטור של פונקציות בסיס $\phi_n(x)$ בקטע $[a,b]$:
    $$f(x) = \sum_{n} c_n \phi_n(x)$$
    כאשר פונקציות הבסיס אורתוגונליות ביחס לפונקציית משקל $w(x) \ge 0$:
    $$\langle \phi_n, \phi_m \rangle = \int_{a}^{b} w(x) \phi_n^*(x) \phi_m(x) \, dx = N_n \delta_{nm}$$
    *(כאשר $N_n = \|\phi_n\|^2$ הוא נורמת הפונקציה בריבוע. אם הבסיס אורתונורמלי, $N_n = 1$)*

    \item \textbf{אלגוריתם חילוץ המקדמים ($c_k$):}
    \begin{enumerate}[leftmargin=*]
        \item \textbf{כפל:} מכפילים את שני אגפי המשוואה בפונקציה $\phi_k^*(x)$ ובפונקציית המשקל $w(x)$.
        \item \textbf{אינטגרציה:} מבצעים אינטגרציה על הקטע $[a,b]$:
        $$\int_{a}^{b} w(x) \phi_k^*(x) f(x) \, dx = \int_{a}^{b} w(x) \phi_k^*(x) \left( \sum_{n} c_n \phi_n(x) \right) dx$$
        \item \textbf{החלפת סדר סכימה ואינטגרציה:}
        $$\int_{a}^{b} w(x) \phi_k^*(x) f(x) \, dx = \sum_{n} c_n \underbrace{\int_{a}^{b} w(x) \phi_k^*(x) \phi_n(x) \, dx}_{N_k \delta_{kn}}$$
        \item \textbf{הפעלת הדלתא של קרונקר ובידוד המקדם:}
        $$c_k = \frac{\langle \phi_k, f \rangle}{\|\phi_k\|^2} = \frac{1}{N_k} \int_{a}^{b} w(x) \phi_k^*(x) f(x) \, dx$$
    \end{enumerate}

    \item \textbf{דוגמאות נפוצות בפיזיקה:}

    \begin{itemize}
        \item \textbf{טור פורייה מורכב ($w(x)=1, \; [-\pi, \pi]$):}
        $$f(x) = \sum_{n=-\infty}^{\infty} c_n e^{i n x} \quad \implies \quad c_k = \frac{1}{2\pi} \int_{-\pi}^{\pi} f(x) e^{-i k x} \, dx$$

        \item \textbf{פיתוח לפי פולינומי לז'נדר ($w(x)=1, \; [-1, 1]$):}
        $$f(x) = \sum_{l=0}^{\infty} A_l P_l(x) \quad \implies \quad A_l = \frac{2l+1}{2} \int_{-1}^{1} f(x) P_l(x) \, dx$$
        *(מכיוון ש-$N_l = \int_{-1}^{1} [P_l(x)]^2 dx = \frac{2}{2l+1}$)*

        \item \textbf{הצלבה מרחבית על ספירה (הרמוניות כדוריות $Y_{lm}(\theta,\phi)$):}
        $f(\theta,\phi) = \sum_{l=0}^{\infty} \sum_{m=-l}^{l} a_{lm} Y_{lm}(\theta,\phi) \Rightarrow a_{lm} = \int_{0}^{2\pi} d\phi \int_{0}^{\pi} \sin\theta d\theta \, Y_{lm}^*(\theta,\phi) f(\theta,\phi)$
    \end{itemize}
\end{itemize}

\section{צפיפות לגראנז'יאן של חלקיק חופשי ומשוואות אוילר-לגראנז'}

    \textbf{חלקיק נקודתי חופשי יחסותי (מכניקה אנליטית):}
    \begin{itemize}
        \item \textbf{לגראנז'יאן לפי זמן עצמי ($\tau$):} $L = -mc \sqrt{-\dot{x}^\mu \dot{x}_\mu}$
        \item \textbf{תוצאת \textenglish{E-L}:} שימור 4-תנע חופשי $\frac{d}{d\tau} p^\mu = 0$ (כאשר $p^\mu = m u^\mu$).
    \end{itemize}


\section{קרינה מטענים מואצים: Bremsstrahlung, Cyclotron ו-Synchrotron}

\begin{itemize}[leftmargin=*]
    \item \textbf{קרינת Bremsstrahlung (קרינת עקירה / בלימה):}
    \begin{itemize}
        \item \textbf{תיאור פיזיקלי:} קרינה הנפלטת כאשר חלקיק טעון $q$ מואט או מוסט במסלולו כתוצאה מאינטראקציה קולונית עם שדה חשמלי של גרעין ($Z e$).
        
        \item \textbf{התפלגות זוויתית של ההספק \textenglish{(CGS)}:}
        כאשר התאוצה מקבילה/אנטי-מקבילה למהירות ($\vec{a} \parallel \vec{v}$, בלימה פרונטלית), קצב פליטת האנרגיה לזווית מרחבית בזמן המטען ניתנת ע"י:
        $$\frac{dP}{d\Omega} = \frac{q^2 a^2 \sin^2\theta}{4\pi c^3 (1-\beta\cos\theta)^5}$$
        *(כאשר $\theta$ היא הזווית בין כיוון התנועה $\vec{\beta}$ לכיוון הצופים $\vec{n}$, ו-$\kappa = 1-\beta\cos\theta$ הוא פקטור הקיטוב היחסותי)*.

        \item \textbf{קירוב יחסותי ($\gamma \gg 1$) ופיתוח לזוויות קטנות ($\theta \ll 1$):}
        נשתמש בקשר היחסותי $\gamma = \frac{E}{m c^2}$ ונגדיר את הפרמטר הקטן $x \equiv \frac{1}{\gamma^2} = \left(\frac{m c^2}{E}\right)^2 \ll 1$.
        לפי קירוב טיילור:
        $$\beta = \sqrt{1 - x} \approx 1 - \frac{1}{2}x$$
        $$\cos\theta \approx 1 - \frac{1}{2}\theta^2, \quad \sin^2\theta \approx \theta^2$$
        מכאן שרכיב המכנה מתפתח ל:
        $$1 - \beta\cos\theta \approx 1 - (1 - \frac{1}{2}x)\left(1 - \frac{1}{2}\theta^2\right) \approx \frac{1}{2}\left(x + \theta^2\right) = \frac{1}{2}\left(\frac{1}{\gamma^2} + \theta^2\right)$$
        ולכן התלות הזוויתית של השטף הינה:
        $$\frac{dP}{d\Omega} \propto \frac{\theta^2}{\left(x + \theta^2\right)^5} = \frac{\theta^2}{\left(\gamma^{-2} + \theta^2\right)^5}$$

        \item \textbf{התנהגות בגבולי הזווית \textenglish{(Relativistic Beaming Limits)}:}
        \begin{enumerate}
            \item \textbf{פנים קונוס הקרינה ($\theta^2 \ll x \iff \theta \ll 1/\gamma$):}
            $$\frac{dP}{d\Omega} \sim \frac{\theta^2}{x^5} = \theta^2 \gamma^{10}$$
            השטף מתאפס במרכז המדויק ($\theta=0$) וגדל ריבועית עם הזווית עד לשיא ב-$\theta_{\max} \approx \frac{1}{2\gamma}$.
            
            \item \textbf{חוץ לקונוס הקרינה ($x \ll \theta^2 \ll 1 \iff \theta \gg 1/\gamma$):}
            $$\frac{dP}{d\Omega} \sim \frac{\theta^2}{\theta^{10}} = \theta^{-8}$$
            דעיכה חזקה מאוד של הקרינה מחוץ לקונוס הצר שזווית פתיחתו היא $\Delta\theta \sim 1/\gamma$.
        \end{enumerate}
    \end{itemize}

    \item \textbf{קרינת Cyclotron (ציקלוטרונית - גבול בלתי-יחסותי):}
    \begin{itemize}
        \item \textbf{תנאים וגבולות:} תנועה מעגלית/ספירלית של חלקיק טעון בשדה מגנטי אחיד $B$ במהירות לא-יחסותית ($\beta \ll 1, \gamma \approx 1$).
        \item \textbf{תכונות המהירות והתאוצה:} $\vec{v} \perp \vec{a}$, תאוצה צנטריפטלית $a = \omega_c v_{\perp}$ כאשר $\omega_c = \frac{|q| B}{m c}$.
        \item \textbf{הספק \textenglish{(Larmor Power)}:} מונוכרומטי בתדירות $\omega_c$, עם הספק כולל $P = \frac{2 q^2 a^2}{3 c^3} = \frac{2 q^4 B^2 v_\perp^2}{3 m^2 c^5}$.
    \end{itemize}



\item \textbf{קרינת Synchrotron (סינכרוטרונית - גבול יחסותי):}
\begin{itemize}
    \item \textbf{תנאים וגבולות:} תנועה בשדה מגנטי $B$ במהירות יחסותית ($\beta \to 1, \; \gamma \gg 1$).
    \item \textbf{תכונות המהירות והתאוצה:} 
    \begin{itemize}
        \item תדירות הסיבוב המסלולית קטנה בגלל הגדלת המסה היחסותית: $\omega_0 = \frac{\omega_c}{\gamma} = \frac{|q| B}{\gamma m c}$.
        \item הקרינה מוצרת לקונוס קדמי בעל זווית פתיחה צרה: $\theta \sim \frac{1}{\gamma}$.
        \item \textbf{הביטוי לתאוצה המאונכת:} כוח לורנץ המגנטי $\vec{F} = \frac{q}{c} (\vec{v} \times \vec{B})$ פועל מאונך לתנועה ($\vec{a} \perp \vec{v} \implies a_\parallel = 0$). מכיוון ש-$\frac{d\vec{p}}{dt} = \gamma m \vec{a}$, גודל התאוצה התלת-ממדית הפיזיקלית הוא:
        $$a_\perp = \frac{|q| B v_\perp}{\gamma m c} = \frac{|q| B \beta_\perp}{\gamma m}$$
    \end{itemize}
    \item \textbf{פיתוח קצר והספק כולל:}
    מביטוי Lienard היחסותי עבור תאוצה מאונכת למהירות ($\vec{a} \perp \vec{v}$):
    $$P = \frac{2 q^2}{3 c^3} \gamma^4 \left( \gamma^2 a_\parallel^2 + a_\perp^2 \right) \xrightarrow{a_\parallel = 0} \frac{2 q^2}{3 c^3} \gamma^4 a_\perp^2$$
    הצבת התאוצה המאונכת $a_\perp = \frac{|q| B \beta_\perp}{\gamma m}$ במשוואה מעניקה:
    $$P_{\text{synch}} = \frac{2 q^2}{3 c^3} \gamma^4 \left( \frac{|q| B \beta_\perp}{\gamma m} \right)^2 = \frac{2 q^4 B^2 \beta_\perp^2 \gamma^2}{3 m^2 c^3}$$
    *(הגורם $\gamma^4$ במונה מצטמצם בחלקו עם $\gamma^2$ במכנה שמקורו בריבוע התאוצה, כך שההספק הסופי פרופורציונלי ל-$\gamma^2$)*.
    \item \textbf{ספקטרום הקרינה:}
    הספקטרום רחב ונפרש עד לתדירות הקריטית:
    $$\omega_c^{\text{synch}} = \frac{3}{2} \gamma^3 \omega_0 \sin\alpha = \frac{3 |q| B \gamma^2}{2 m c} \sin\alpha$$
    *(כאשר $\alpha$ היא זווית ה-Pitch בין המהירות $\vec{v}$ לשדה המגנטי $\vec{B}$, כך ש-$v_\perp = v \sin\alpha$)*.
\end{itemize}


\end{itemize}

\section*{הטנזור הדואלי \textenglish{(CGS and Mostly Minus)}:}
    {\scriptsize$
    \tilde{F}^{\mu\nu} = \begin{pmatrix}
    0 & -B_x & -B_y & -B_z \\
    B_x & 0 & E_z & -E_y \\
    B_y & -E_z & 0 & E_x \\
    B_z & E_y & -E_x & 0
    \end{pmatrix}, \tilde{F}_{\mu\nu} = \begin{pmatrix}
    0 & B_x & B_y & B_z \\
    -B_x & 0 & E_z & -E_y \\
    -B_y & -E_z & 0 & E_x \\
    -B_z & E_y & -E_x & 0
    \end{pmatrix}$} \\
    *(כאשר $\epsilon^{0123} = +1$)*.

\section{זיהוי וחילוץ המולטיפול המוביל בקרינה מתוך הפוטנציאל הווקטורי}

\begin{itemize}[leftmargin=*]
    \item \textbf{הבסיס הפיזיקלי והפיתוח בשדה הרחוק ($R \to \infty$):}
    בשדה הרחוק ($r \gg r'$ ו-$r \gg \lambda$), הפוטנציאל הווקטורי המעוכב $\vec{A}(\vec{r}, t)$ מפותח כטור חזקות של גודל המקור ביחס למרחק ($r'/r$) וביחס לאורך הגל ($r'/\lambda \sim v/c$):
    $$\vec{A}(\vec{r}, t) = \frac{1}{c r} \int \vec{J}\left(\vec{r}', t - \frac{r}{c} + \frac{\hat{n} \cdot \vec{r}'}{c}\right) d^3r'$$
    פיתוח טיילור של הזמן המעוכב סביב $t_0 = t - \frac{r}{c}$ מניב:
    $$\vec{A}(\vec{r}, t) \approx \frac{1}{c r} \int \vec{J}(\vec{r}', t_0) d^3r' + \frac{1}{c^2 r} \frac{\partial}{\partial t} \int (\hat{n} \cdot \vec{r}') \vec{J}(\vec{r}', t_0) d^3r' + \dots$$

    \item \textbf{כללי זיהוי המולטיפול המוביל \textenglish{(Leading Order)}:}
    המולטיפול המוביל הוא **האיבר הראשון בטור שאינו מתאפס**:
    \begin{enumerate}
        \item \textbf{דיפול חשמלי \textenglish{(Electric Dipole - E1)}:} האיבר הראשון ($r'^0$). מתוך זהות הרציפות ($\nabla \cdot \vec{J} = -\dot{\rho}$):
        $$\int \vec{J} d^3r' = \dot{\vec{p}} \quad \implies \quad \vec{A}_{\text{E1}} = \frac{\dot{\vec{p}}(t_0)}{c r}$$
        *אם $\dot{\vec{p}} \neq 0$, הקרינה היא מסדר דיפול חשמלי \textenglish{(E1)}.*

        \item \textbf{איבר הסדר השני ($r'^1$):} אם $\dot{\vec{p}} = 0$, מפרקים את האיבר השני לחלק אנטי-סימטרי (דיפול מגנטי \textenglish{M1}) ולחלק סימטרי (קוודרופול חשמלי \textenglish{E2}):
        $$\vec{A}^{(2)} = \underbrace{\frac{\dot{\vec{m}}(t_0) \times \hat{n}}{r}}_{\text{\textenglish{M1 (Magnetic Dipole)}}} + \underbrace{\frac{\ddot{\mathbf{Q}}(t_0) \cdot \hat{n}}{6 c^2 r}}_{\text{\textenglish{E2 (Electric Quadrupole)}}}$$
    \end{enumerate}

    \item \textbf{דוגמה חישובית: טבעת זרם רוטטת / משתנה בזמן}
    \begin{itemize}
        \item \textbf{תיאור המערכת:} טבעת דקה ברדיוס $R$ במישור $xy$, שזורם בה זרם אחיד משתנה בזמן $I(t) = I_0 \cos(\omega t) \hat{\phi}$.
        \item \textbf{שלב 1 — בדיקת דיפול חשמלי \textenglish{(E1)}:}
        צפיפות המטען מתאפסת בכל המרחב ($\rho = 0$), ולכן המומנט הדיפולי החשמלי מתאפס תמיד:
        $$\vec{p} = \int \rho \vec{r}' d^3r' = 0 \quad \implies \quad \vec{A}_{\text{E1}} = 0$$
        *(אין קרינת דיפול חשמלי)*.

        \item \textbf{שלב 2 — חישוב האיבר הבא (\textenglish{M1} לעומת \textenglish{E2}):}
        נחשב את המומנט המגנטי של הטבעת:
        $$\vec{m}(t) = \frac{1}{2c} \int (\vec{r}' \times \vec{J}) d^3r' = \frac{I(t)}{c} \cdot (\pi R^2) \hat{z} = \frac{\pi R^2 I_0 \cos(\omega t)}{c} \hat{z}$$
        מכיוון ש-$\dot{\vec{m}}(t) = -\frac{\pi R^2 I_0 \omega \sin(\omega t)}{c} \hat{z} \neq 0$, הפוטנציאל הווקטורי המוביל בשדה הרחוק נתון ע"י:
        $$\vec{A}_{\text{rad}}(\vec{r}, t) \approx \frac{\dot{\vec{m}}(t_0) \times \hat{n}}{r} = -\frac{\pi R^2 I_0 \omega \sin(\omega t_0)}{c r} (\hat{z} \times \hat{n}), \ (\hat{n} = \hat{r})$$
        \item \textbf{מסקנה:} המולטיפול המוביל בקרינה עבור טבעת זרם זו הוא **דיפול מגנטי \textenglish{(M1)}**.
    \end{itemize}
\end{itemize}

\end{multicols*}
\end{document}